\documentclass[11pt,a4paper]{article}
% Preâmbulo comum. Não duplicar pacotes nos módulos.
\usepackage[utf8]{inputenc}
\usepackage[T1]{fontenc}
\usepackage[brazil]{babel}
\usepackage{amsmath,amssymb,amsthm,mathtools}
\usepackage{geometry}
\geometry{a4paper,margin=2.3cm}
\usepackage{microtype}
\usepackage{enumitem}
\usepackage{graphicx}
\usepackage{xcolor}
\usepackage{booktabs}
\usepackage{hyperref}
\hypersetup{colorlinks=true,linkcolor=blue!40!black,urlcolor=blue!50!black}

\setcounter{tocdepth}{2}
\numberwithin{equation}{section}

\theoremstyle{plain}
\newtheorem{teorema}{Teorema}[section]
\newtheorem{proposicao}[teorema]{Proposição}
\newtheorem{lema}[teorema]{Lema}
\newtheorem{corolario}[teorema]{Corolário}

\theoremstyle{definition}
\newtheorem{definicao}[teorema]{Definição}
\newtheorem{exemplo}[teorema]{Exemplo}
\newtheorem{exercicio}{Exercício}[section]

\theoremstyle{remark}
\newtheorem{observacao}[teorema]{Observação}
\newtheorem{contraexemplo}[teorema]{Contraexemplo}
\newtheorem{intuicao}[teorema]{Intuição}

\newcommand{\R}{\mathbb{R}}
\newcommand{\N}{\mathbb{N}}
\newcommand{\Z}{\mathbb{Z}}
\newcommand{\Q}{\mathbb{Q}}
\DeclareRobustCommand{\eps}{\varepsilon}
\DeclareRobustCommand{\vet}[1]{\vec{#1}}
\DeclareRobustCommand{\dif}{\mathrm{d}}
\DeclareRobustCommand{\abs}[1]{\left\lvert #1 \right\rvert}
\DeclareRobustCommand{\norm}[1]{\left\lVert #1 \right\rVert}
\DeclareMathOperator{\grad}{grad}
\DeclareMathOperator{\Div}{div}
\DeclareMathOperator{\rot}{rot}

\newcommand{\base}[1]{\par\medskip\noindent\textbf{Base para relembrar.} #1\par\medskip}
\newcommand{\cuidado}[1]{\par\medskip\noindent\textbf{Quando a hipótese falha.} #1\par\medskip}

\newcommand{\figura}[2]{%
  \begin{figure}[htbp]
  \centering
  \IfFileExists{#1}{%
    \includegraphics[width=0.84\textwidth]{#1}%
  }{%
    \fbox{\parbox{0.8\textwidth}{\centering Figura ainda não gerada:\\ \texttt{#1}}}%
  }
  \caption{#2}
  \end{figure}%
}

\newenvironment{esqueleto}{%
  \par\medskip\noindent\textbf{Esqueleto da demonstração.}\par\nopagebreak
  \begin{enumerate}[leftmargin=1.4em]
}{\end{enumerate}\medskip}

\title{Caderno de prova: séries e Taylor}
\author{}
\date{}
\begin{document}
\maketitle
\tableofcontents

Os enunciados não trazem roteiro. O espaço em branco é para a prova.

\section{Enunciados}

\begin{exercicio}
Seja \(r\in\R\). Prove que, para todo \(N\in\N\),
\[
\sum_{k=0}^{N-1} r^k=
\begin{cases}
N, & \text{se }r=1,\\[0.4em]
\dfrac{1-r^N}{1-r}, & \text{se }r\neq 1.
\end{cases}
\]
Prove em seguida que a série \(\sum_{k=0}^\infty r^k\) converge se e somente se \(\abs{r}<1\), e que nesse caso a soma é \(1/(1-r)\).
\end{exercicio}

\vspace{7cm}

\begin{exercicio}
Prove que, se \(\sum_{n=1}^\infty a_n\) converge, então \(a_n\to 0\). Depois prove que \(\sum_{n=1}^\infty 1/n\) diverge, e conclua que a recíproca da primeira afirmação é falsa.
\end{exercicio}

\vspace{7cm}

\begin{exercicio}
Suponha que existe
\[
L=\lim_{n\to\infty}\abs{\frac{a_{n+1}}{a_n}}
\]
e que \(L<1\). Prove que \(\sum a_n\) converge absolutamente.
\end{exercicio}

\vspace{7cm}

\begin{exercicio}
Seja \((b_n)\) uma sequência decrescente, com \(b_n\ge 0\) e \(b_n\to 0\). Prove que
\[
\sum_{n=1}^\infty (-1)^{n+1} b_n
\]
converge. Se \(S\) é a soma e \(s_N\) é a soma parcial, prove que \(\abs{S-s_N}\le b_{N+1}\).
\end{exercicio}

\vspace{7cm}

\begin{exercicio}
Suponha que \(f''\) existe num intervalo aberto que contém o segmento de extremos \(a\) e \(x\), com \(x\neq a\). Seja \(T_1(x)=f(a)+f'(a)(x-a)\). Prove que existe \(c\) estritamente entre \(a\) e \(x\) tal que
\[
f(x)-T_1(x)=\frac{f''(c)}{2}(x-a)^2.
\]
Pode usar o teorema de Rolle: se uma função é contínua no segmento fechado, derivável no aberto e toma o mesmo valor nos dois extremos, então a derivada se anula em algum ponto estritamente entre eles.
\end{exercicio}

\vspace{7cm}

\begin{exercicio}
Prove que, para todo \(x\in\R\),
\[
e^x=\sum_{n=0}^\infty\frac{x^n}{n!}.
\]
Pode usar, como teorema já demonstrado, o resto de Lagrange de ordem \(n\): se \(f^{(n+1)}\) existe num intervalo aberto que contém o segmento de extremos \(0\) e \(x\), então existe \(c\) estritamente entre \(0\) e \(x\) com
\[
f(x)-\sum_{k=0}^n\frac{f^{(k)}(0)}{k!}x^k=\frac{f^{(n+1)}(c)}{(n+1)!}x^{n+1},
\]
e o resto é zero se \(x=0\).
\end{exercicio}

\vspace{7cm}

\end{document}
